% Compile with XeLaTeX.
\documentclass[11pt,a4paper]{ctexart}
\usepackage[margin=2.3cm]{geometry}
\usepackage{amsmath,amssymb}
\usepackage{enumitem}
\usepackage{hyperref}
\hypersetup{colorlinks=true,linkcolor=blue,urlcolor=blue}
\setlength{\parindent}{0pt}
\setlength{\parskip}{0.5em}
\providecommand{\tightlist}{\setlength{\itemsep}{0pt}\setlength{\parskip}{0pt}}
\setcounter{secnumdepth}{0}
\begin{document}
\hypertarget{gre-mathematics-ux9519ux9898ux4e0eux4e0dux786eux5b9aux9898practice-1}{%
\section{GRE Mathematics 错题与不确定题：Practice
1}\label{gre-mathematics-ux9519ux9898ux4e0eux4e0dux786eux5b9aux9898practice-1}}

本次记录 10
道错题：24、27、43、44、50、53、60、62、64、65；另收录答对但不确定的
36、37。题号和选项对应 Practice
1。错因依据作答者自述；错题不自动等于知识漏洞，猜对也不等于掌握。

\hypertarget{ux9f50ux6b21ux65b9ux7a0bux7ec4ux7684ux89e3ux7a7aux95f4e}{%
\subsection{24.
齐次方程组的解空间：E}\label{ux9f50ux6b21ux65b9ux7a0bux7ec4ux7684ux89e3ux7a7aux95f4e}}

\textbf{原题（Original problem）}\par
Consider the system of linear equations
\[\begin{aligned}w+3x+2y+2z&=0,\\ w+4x+y&=0,\\ 3w+5x+10y+14z&=0,\\ 2w+5x+5y+6z&=0.\end{aligned}\]
with solutions of the form $(w,x,y,z)$, where $w$, $x$, $y$, and $z$ are real. Which of the following statements is FALSE?
\begin{enumerate}[label=(\Alph*)]
\item The system is consistent. \item The system has infinitely many solutions. \item The sum of any two solutions is a solution. \item $(-5,1,1,0)$ is a solution. \item Every solution is a scalar multiple of $(-5,1,1,0)$.
\end{enumerate}
\textbf{词语：}consistent：有解（has at least one solution）；scalar multiple：标量倍数；FALSE：错误的。

\textbf{解答与分析}\par


\textbf{错因：}误解 consistent，且未验证解空间维数。consistent
意为至少有一个解；齐次方程组一定有零解，因此必 consistent。

第二式给 \(w=-4x-y\)，代入第一式得 \(x=y+2z\)，因此

\[ (w,x,y,z)=y(-5,1,1,0)+z(-8,2,0,1). \]

另外两式也满足。两个向量线性无关，所以解空间维数为
2，并非所有解都是第一个向量的倍数，E 错。A 正确；B 有无穷多解；C
因线性性成立；D 可直接代入。一般而言齐次系统 \(Ax=0\)
的解构成子空间，维数为未知数个数减去
\(\operatorname{rank}A\)；非齐次系统有解才称
consistent，其解集一般是一个仿射空间。

\hypertarget{ux590dux6570ux7684ux5e45ux89d2ux4e0eux6a21d}{%
\subsection{27.
复数的幅角与模：D}\label{ux590dux6570ux7684ux5e45ux89d2ux4e0eux6a21d}}

\textbf{原题（Original problem）}\par
$(1+i)^{10}=$
\begin{enumerate}[label=(\Alph*),nosep]
\item $1$ \item $i$ \item $32$ \item $32i$ \item $32(i+1)$
\end{enumerate}
\textbf{词语：}modulus：复数的模；argument：复数的幅角。

\textbf{解答与分析}\par


\textbf{错因：}只考虑幅角，漏掉模长。

\[ (1+i)^{10}=\bigl(\sqrt2\,e^{i\pi/4}\bigr)^{10}=32e^{5\pi i/2}=32i. \]

也可用 \((1+i)^2=2i\)，得到
\((2i)^5=32i\)。复数乘方时，模取幂、幅角乘指数；两者缺一不可。

\hypertarget{ux5185ux90e8ux95edux5305ux548cux8fdeux901acux7b54ux5bf9ux4f46ux4e0dux786eux5b9a}{%
\subsection{36.
内部、闭包和连通：C（答对但不确定）}\label{ux5185ux90e8ux95edux5305ux548cux8fdeux901acux7b54ux5bf9ux4f46ux4e0dux786eux5b9a}}

\textbf{原题（Original problem）}\par
Suppose $S$ is a nonempty subset of $\mathbb R$. Which of the following is necessarily true?
\begin{enumerate}[label=(\Alph*)]
\item For each $s,t\in S$, there exists a continuous function $f$ mapping $[0,1]$ into $S$ with $f(0)=s$ and $f(1)=t$.
\item For each $u\notin S$, there exists an open subset $U$ of $\mathbb R$ such that $u\in U$ and $U\cap S=\varnothing$.
\item $\{v\in S:\text{there exists an open subset }V\text{ of }\mathbb R\text{ with }v\in V\subseteq S\}$ is an open subset of $\mathbb R$.
\item $\{w\notin S:\text{there exists an open subset }W\text{ of }\mathbb R\text{ with }w\in W\text{ and }W\cap S=\varnothing\}$ is a closed subset of $\mathbb R$.
\item $S$ is the intersection of all closed subsets of $\mathbb R$ that contain $S$.
\end{enumerate}
\textbf{词语：}necessarily true：必然成立；open/closed subset：开/闭子集；intersection：交集。

\textbf{解答与分析}\par


C 所述集合恰是内部 \(S^\circ\)：所有存在开邻域包含于 \(S\)
的点。它是所有包含于 \(S\) 的开集之并，因而开。

\begin{itemize}
\tightlist
\item
  A 要求任意两点之间有包含于 \(S\) 的连续路径；取 \(S=\{0,1\}\)
  即不成立，因为连续函数在区间上的像是连通的。
\item
  B 要求每个不属于 \(S\) 的点有不碰 \(S\) 的开邻域；取
  \(S=(0,1)\)、\(u=0\)，不成立。
\item
  D 所述集合是 \(\mathbb R\setminus\overline S\)，总是开集，不一定闭。取
  \(S=[0,1]\)。
\item
  E 中所有包含 \(S\) 的闭集之交是闭包 \(\overline S\)，不一定是
  \(S\)；取 \(S=(0,1)\)。
\end{itemize}

\textbf{速记：}内部是最大开子集，闭包是最小闭超集，边界为
\(\overline S\setminus S^\circ\)。集合开不代表不闭，闭不代表不开。

\hypertarget{ux5e42ux7b49ux53d8ux6362cux7b54ux5bf9ux4f46ux4e0dux786eux5b9a}{%
\subsection{37.
幂等变换：C（答对但不确定）}\label{ux5e42ux7b49ux53d8ux6362cux7b54ux5bf9ux4f46ux4e0dux786eux5b9a}}

\textbf{原题（Original problem）}\par
Let $V$ be a finite-dimensional real vector space and let $P$ be a linear transformation of $V$ such that $P^2=P$. Which of the following must be true?
\begin{enumerate}[label=\Roman*.]
\item $P$ is invertible. \item $P$ is diagonalizable. \item $P$ is either the identity transformation or the zero transformation.
\end{enumerate}
(A) None \quad (B) I only \quad (C) II only \quad (D) III only \quad (E) II and III.

\textbf{词语：}invertible：可逆；diagonalizable：可对角化；identity transformation：恒等变换。

\textbf{解答与分析}\par


\textbf{基本概念。}像空间与核空间分别定义为
\[
\operatorname{im}P=\{P(v):v\in V\},\qquad
\ker P=\{v\in V:P(v)=0\}.
\]

\textbf{第一步：像空间中的向量经过 \(P\) 后不变。}任取 \(u\in\operatorname{im}P\)，由定义存在 \(v\in V\) 使 \(u=P(v)\)。因此
\[
P(u)=P(P(v))=P^2(v)=P(v)=u.
\]
所以 \(P\) 在像空间上是恒等变换；其中的每个非零向量都是特征值 \(1\) 的特征向量。

\textbf{第二步：核空间中的向量经过 \(P\) 后变为零。}若 \(w\in\ker P\)，则按核的定义 \(P(w)=0\)。因此核空间中的每个非零向量都是特征值 \(0\) 的特征向量。

\textbf{第三步：证明 \(V=\operatorname{im}P\oplus\ker P\)。}任取 \(v\in V\)，写成
\[
v=P(v)+\bigl(v-P(v)\bigr).
\]
第一项 \(P(v)\in\operatorname{im}P\)。对第二项作用 \(P\)，由线性性与 \(P^2=P\)，
\[
P\bigl(v-P(v)\bigr)=P(v)-P^2(v)=0,
\]
故 \(v-P(v)\in\ker P\)。这证明每个向量都是像空间中的向量与核空间中的向量之和，即 \(V=\operatorname{im}P+\ker P\)。

这个分解还是唯一的：若 \(u\in\operatorname{im}P\cap\ker P\)，由第一步有 \(P(u)=u\)，由核的定义又有 \(P(u)=0\)，故 \(u=0\)。因此两个子空间交集仅含零向量，得到直和
\[
V=\operatorname{im}P\oplus\ker P.
\]

\textbf{第四步：选基并写出矩阵。}取像空间的一组基 \(u_1,\ldots,u_r\)，以及核空间的一组基 \(w_1,\ldots,w_s\)。上述直和分解保证
\[
\mathcal B=(u_1,\ldots,u_r,w_1,\ldots,w_s)
\]
是 \(V\) 的一组基。对这些基向量，
\[
P(u_i)=u_i\quad(1\le i\le r),\qquad P(w_j)=0\quad(1\le j\le s).
\]
矩阵的每一列是相应基向量的像在基 \(\mathcal B\) 下的坐标，因此
\[
[P]_{\mathcal B}=
\begin{pmatrix}I_r&0\\0&0\end{pmatrix}
=\operatorname{diag}(\underbrace{1,\ldots,1}_{r},
\underbrace{0,\ldots,0}_{s}).
\]
这是一個对角矩阵。于是 \(P\) 有一组由特征向量组成的基，故 \(P\) 可对角化，II 成立。若像空间或核空间是零空间，则相应的一组基为空，上述结论仍成立。

\textbf{具体例子。}在 \(\mathbb R^2\) 上令 \(P(x,y)=(x,0)\)。则 \(P^2=P\)，像空间为 \(x\) 轴 \(\operatorname{span}\{(1,0)\}\)，核空间为 \(y\) 轴 \(\operatorname{span}\{(0,1)\}\)。在由这两个向量组成的基下，矩阵为
\[
\begin{pmatrix}1&0\\0&0\end{pmatrix}.
\]
I 不一定成立：零变换在非零空间上不可逆。III 不一定成立：这个例子既非恒等变换也非零变换。因此答案为 C。

\textbf{易错点。}仅证明特征值属于 \(\{0,1\}\) 本身不足以证明可对角化；还需要证明存在足够多的线性无关特征向量。上述直和分解给出了完整特征向量基。也可使用最小多项式整除无重根多项式 \(t(t-1)\) 来证明。

\hypertarget{ux53c2ux6570ux66f2ux7ebfux4e8cux9636ux5bfcux6570a}{%
\subsection{43.
参数曲线二阶导数：A}\label{ux53c2ux6570ux66f2ux7ebfux4e8cux9636ux5bfcux6570a}}

\textbf{原题（Original problem）}\par
A curve in the $xy$-plane is given parametrically by
\[x=t^2+2t,\qquad y=3t^4+4t^3\]
for all $t>0$. The value of $d^2y/dx^2$ at the point $(8,80)$ is

(A) $4$ \quad (B) $24$ \quad (C) $32$ \quad (D) $96$ \quad (E) $192$.

\textbf{词语：}given parametrically：以参数方程给出；second derivative：二阶导数。

\textbf{解答与分析}\par


\textbf{错因：}求解过快。\(x=t^2+2t,\ y=3t^4+4t^3\)，由 \(t>0\) 和
\(x=8\) 得 \(t=2\)。先求

\[\frac{dy}{dx}=\frac{12t^3+12t^2}{2t+2}=6t^2,\qquad \frac{d^2y}{dx^2}=\frac{d(6t^2)/dt}{dx/dt}=\frac{12t}{2t+2}.\]

代入 \(t=2\) 得 4。每次对 \(x\) 求导，都要除以 \(dx/dt\)；不是直接对
\(t\) 求两次导数。

\hypertarget{ux4e00ux9636ux5faeux5206ux65b9ux7a0bb}{%
\subsection{44.
一阶微分方程：B}\label{ux4e00ux9636ux5faeux5206ux65b9ux7a0bb}}

\textbf{原题（Original problem）}\par
\[y'+xy=x,\qquad y(0)=-1.\]
If $y$ is a real-valued function defined on the real line and satisfying the initial value problem above, then $\displaystyle\lim_{x\to-\infty}y(x)=$

(A) $0$ \quad (B) $1$ \quad (C) $-1$ \quad (D) $\infty$ \quad (E) $-\infty$.

\textbf{词语：}initial value problem：初值问题；real-valued：取实值的；satisfying：满足。

\textbf{解答与分析}\par


\textbf{错因：}符号未代回原式核验。\(y'+xy=x,\ y(0)=-1\)。令
\(u=y-1\)，则 \(u'=-xu\)，所以

\[y=1+Ce^{-x^2/2}=1-2e^{-x^2/2},\qquad \lim_{x\to-\infty}y(x)=1.\]

检验：\(y'=2xe^{-x^2/2}\)，故 \(y'+xy=x\)，且 \(y(0)=-1\)。分离变量时
\(\int -x\,dx=-x^2/2\)；绝对值和任意常数应处理正确。

\hypertarget{ux7406ux60f3ux7684ux548cux9010ux9879ux4e58ux79efux4e0eux4ea4d}{%
\subsection{50.
理想的和、逐项乘积与交：D}\label{ux7406ux60f3ux7684ux548cux9010ux9879ux4e58ux79efux4e0eux4ea4d}}

\textbf{原题（Original problem）}\par
Let $R$ be a ring and let $U$ and $V$ be (two-sided) ideals of $R$. Which of the following must also be ideals of $R$?
\begin{enumerate}[label=\Roman*.]
\item $U+V=\{u+v:u\in U\text{ and }v\in V\}$.
\item $U\cdot V=\{uv:u\in U\text{ and }v\in V\}$.
\item $U\cap V$.
\end{enumerate}
(A) II only \quad (B) III only \quad (C) I and II only \quad (D) I and III only \quad (E) I, II, and III.

\textbf{词语：}ring：环；(two-sided) ideal：（双侧）理想；must also be：也必定是。注意 II 在题目中明确指单个乘积的集合。

\textbf{解答与分析}\par


\textbf{错因：}理想的闭合条件及理想乘积的定义不熟。双侧理想 \(I\)
是加法子群，且对 \(r\in R,a\in I\)，有 \(ra,ar\in I\)。

I：\(U+V=\{u+v\}\) 是理想；加法与左右乘吸收分别由 \(U,V\)
的性质得到。III：\(U\cap V\) 是理想，因为相关运算同时留在 \(U\) 和
\(V\)。

II：题目只允许\textbf{单个乘积}构成的集合
\(\{uv:u\in U,v\in V\}\)，不一定对加法封闭。具体取
\(R=\mathbb R[x,y]\)、\(U=V=(x,y)\)。\(x^2,y^2\) 都是单个乘积，但
\(x^2+y^2\) 不是两个属于 \((x,y)\)
的多项式的乘积：若能分解，两个因子的最低非零齐次部分均为实线性式，它们的乘积应为
\(x^2+y^2\)，然而后者不能分解为两个实线性式。故 II 不成立。

\textbf{相关高频知识：}

\begin{enumerate}
\def\labelenumi{\arabic{enumi}.}
\tightlist
\item
  通常的理想乘积 \(UV\) 定义为有限和
  \(\{\sum_{i=1}^m u_iv_i\}\)，不是单个乘积集合；它总是双侧理想，且
  \(UV\subseteq U\cap V\)。理想的并一般不是理想；两个理想的并是理想当且仅当其中一个包含另一个。
\item
  环同态的核是理想；商环 \(R/I\) 需要 \(I\)
  是双侧理想。交换含幺环中，真理想 \(I\) 极大当且仅当 \(R/I\)
  是域；\(I\) 素当且仅当 \(R/I\) 是整环。域的理想只有 \(\{0\}\)
  和整个域。
\item
  \(\mathbb Z\) 的每个理想都是
  \(n\mathbb Z\)；\(a\mathbb Z+b\mathbb Z=\gcd(a,b)\mathbb Z\)，交为
  \(\operatorname{lcm}(a,b)\mathbb Z\)。\(\mathbb Z/n\mathbb Z\)
  是域当且仅当 \(n\) 为素数。
\item
  群同态的核是正规子群；有限群中元素的阶整除群的阶（拉格朗日定理），但阶整除不意味着一定存在该阶元素。循环群的子群均循环；阶为
  \(n\) 的循环群中 \(g^k\) 的阶为 \(n/\gcd(n,k)\)。
\end{enumerate}

\hypertarget{ux53d8ux4e0aux9650ux79efux5206ux4e0eux5149ux6ed1ux6027a}{%
\subsection{53.
变上限积分与光滑性：A}\label{ux53d8ux4e0aux9650ux79efux5206ux4e0eux5149ux6ed1ux6027a}}

\textbf{原题（Original problem）}\par
Let $f$ and $g$ be functions of a real variable such that $g(x)=\int_0^x f(y)(y-x)\,dy$ for all $x$. If $g$ is three times continuously differentiable, what is the greatest integer $n$ for which $f$ must be $n$ times continuously differentiable?

(A) $1$ \quad (B) $2$ \quad (C) $3$ \quad (D) $4$ \quad (E) $5$.

\textbf{词语：}continuously differentiable：连续可微；greatest integer：最大整数；must be：必然满足。

\textbf{解答与分析}\par


题设 \(g(x)=\int_0^x f(y)(y-x)\,dy\)，给定 \(g\in C^3\)，求必须有
\(f\in C^n\) 的最大 \(n\)。解题时要注意方向：从 \(g\)
的正则性\textbf{反推} \(f\)。

在通常的连续被积函数假设下，莱布尼茨公式给出

\[g'(x)=-\int_0^x f(y)\,dy,\qquad g''(x)=-f(x),\qquad g'''(x)=-f'(x).\]

因此 \(f=-g''\in C^1\)。不一定是 \(C^2\)：取 \(f(x)=|x|^{3/2}\)，它属于
\(C^1\) 但在 0 不属于 \(C^2\)；由上式构造的 \(g\) 属于
\(C^3\)，故不能保证更多。

关于题设的细节：原题未明确写 \(f\) 连续。若只按通常 Riemann
积分理解，改变 \(f\) 在单个点的值不改变任何积分，却可破坏 \(f\)
的连续性；于是严格字面表述有缺口，A 需要补充 \(f\) 连续等假设。若按广义
Lebesgue 积分，更只能先得到几乎处处的结论。ETS 预期答案 A
基于通常的连续性假设，不应把这一隐含条件忽略。

你的直觉抓住了变上限积分带来的求导关系；这里具体是两次求导后恢复
\(-f\)，所以 \(g\) 的三次连续可微对应 \(f\) 的一次连续可微。

\hypertarget{ux91cfux8bcdux548cux65e0ux7a77ux8fdcux5904ux6781ux9650d}{%
\subsection{60.
量词和无穷远处极限：D}\label{ux91cfux8bcdux548cux65e0ux7a77ux8fdcux5904ux6781ux9650d}}

\textbf{原题（Original problem）}\par
A real-valued function $f$ defined on $\mathbb R$ has the following property.

For every positive number $\epsilon$, there exists a positive number $\delta$ such that $|f(x)-f(1)|\ge\epsilon$ whenever $|x-1|\ge\delta$.

This property is equivalent to which of the following statements about $f$?
\begin{enumerate}[label=(\Alph*)]
\item $f$ is continuous at $x=1$. \item $f$ is discontinuous at $x=1$. \item $f$ is unbounded. \item $\displaystyle\lim_{|x|\to\infty}|f(x)|=\infty$. \item $\displaystyle\int_0^\infty |f(x)|\,dx=\infty$.
\end{enumerate}
\textbf{词语：}whenever：只要；equivalent to：等价于；unbounded：无界。

\textbf{解答与分析}\par


题设是对每个 \(\varepsilon>0\)，存在 \(\delta>0\)，只要
\textbf{\(|x-1|\ge\delta\)}，就有
\textbf{\(|f(x)-f(1)|\ge\varepsilon\)}。它控制的是远离 1
的点，输出也远离 \(f(1)\)，不是连续性。连续性要求
\(|x-1|<\delta\Rightarrow|f(x)-f(1)|<\varepsilon\)。

令 \(c=f(1)\)。题设等价于 \(|f(x)-c|\to\infty\)，当
\(|x|\to\infty\)：因为 \(|x|\) 与 \(|x-1|\) 同时趋向无穷。再由

\[\bigl||f(x)|-|f(x)-c|\bigr|\le |c|\]

可知它等价于 \(|f(x)|\to\infty\)，即 D。形式上，如要证明 D，由题设选阈值
\(M+|c|+1\)，即可得到 \(|f(x)|>M\)；反向同理。

A 不等价：\(f(x)=x\)
连续且满足题设，但常数函数也连续，却不满足题设；另外令
\(f(1)=0\)、\(f(x)=x\)（\(x\ne1\)），满足 D 但在 1 不连续。C
仅要求无界，如 \(\sin(x^2)x\) 无界却不趋于无穷。不要只盯
\(\varepsilon,\delta\)，先看两个不等号的方向与量词。

\hypertarget{ux5750ux6807ux4e2dux81f3ux5c11ux4e00ux4e2aux65e0ux7406c}{%
\subsection{62.
坐标中至少一个无理：C}\label{ux5750ux6807ux4e2dux81f3ux5c11ux4e00ux4e2aux65e0ux7406c}}

\textbf{原题（Original problem）}\par
Let $S$ be the subset of $\mathbb R^2$ consisting of all points $(x,y)$ in the unit square $[0,1]\times[0,1]$ for which $x$ or $y$, or both, are irrational. With respect to the standard topology on $\mathbb R^2$, $S$ is
\begin{enumerate}[label=(\Alph*),nosep]
\item closed \item open \item connected \item totally disconnected \item compact
\end{enumerate}
\textbf{词语：}or both：或者两者都是，即至少一个；irrational：无理数；connected：连通；totally disconnected：全不连通；compact：紧。

\textbf{解答与分析}\par


\textbf{错因：}漏读 ``x or y, or both''。设
\(S=[0,1]^2\setminus(\mathbb Q\times\mathbb Q)\)。从任意
\(p=(a,b)\in S\) 出发，若 \(a\) 无理，先沿竖线到
\((a,\beta)\)，其中任选无理 \(\beta\in(0,1)\)；若 \(a\) 有理，则 \(b\)
无理，先水平走到 \((\alpha,b)\)（\(\alpha\) 无理），再竖直走到
\((\alpha,\beta)\)。每段至少有一个固定的无理坐标。任意两点均可连到同一个
\((\alpha,\beta)\)，所以 \(S\) 路径连通，从而连通。

\(S\) 在正方形中稠密，但不包含有理坐标点，因此在 \(\mathbb R^2\)
中不闭、不紧；也不是开集，因为每个邻域都含有两个坐标均有理的点。它包含线段，所以不是全不连通。注意开闭均相对于题目指定的
\(\mathbb R^2\)。

\hypertarget{ux6563ux5ea6ux5b9aux7406ux901aux91cfux53caux573aux8bbae}{%
\subsection{64.
散度定理、通量及场论：E}\label{ux6563ux5ea6ux5b9aux7406ux901aux91cfux53caux573aux8bbae}}

\textbf{原题（Original problem）}\par
What is the value of the flux of the vector field $\mathbf F$, defined on $\mathbb R^3$ by $\mathbf F(x,y,z)=x\mathbf i+y\mathbf j+z\mathbf k$, through the surface $z=\sqrt{1-x^2-y^2}$ oriented with upward-pointing normal vector field? (Note: $\mathbf i$, $\mathbf j$, and $\mathbf k$ are the standard basis vectors in $\mathbb R^3$.)

(A) $0$ \quad (B) $2\pi/3$ \quad (C) $\pi$ \quad (D) $4\pi/3$ \quad (E) $2\pi$.

\textbf{词语：}flux：通量；vector field：向量场；upward-pointing normal：指向上方的法向量；oriented：指定方向的。

\textbf{解答与分析}\par


题目是
\(z=\sqrt{1-x^2-y^2}\)，即单位球的上半球，而不是抛物面。\(F=(x,y,z)\)，散度为
3。用 \(z=0\) 的圆盘封闭半球，半球体积为 \(2\pi/3\)；圆盘外法向为
\(-\mathbf k\)，由于 \(z=0\)，底面通量为 0。因此

\[\iint_S F\cdot n\,dS=3\cdot\frac{2\pi}{3}=2\pi.\]

另一种快速方法：单位球上外法向 \(n=(x,y,z)=F\)，所以
\(F\cdot n=1\)，通量就是上半球面积 \(2\pi\)。你的红笔计算 \(4\pi\)
是整个单位球的总通量；本题只取上半球，因此减半。

\textbf{通用方法：}图形曲面 \(z=h(x,y)\) 上向面积向量是
\((-h_x,-h_y,1)\,dx\,dy\)，因此通量
\(\iint_D[-P h_x-Q h_y+R]dxdy\)；下向取负。若能加盖形成闭曲面，先算
\(\iiint_\Omega\operatorname{div}F\,dV\)，再减去盖面的通量。换成其他
\(F=(P,Q,R)\) 时，先算
\(\operatorname{div}F=P_x+Q_y+R_z\)，不要把散度默认为 3。

\textbf{三个核心公式：}

\begin{itemize}
\tightlist
\item
  Green：正向平面简单闭曲线
  \(C=\partial D\)，\(\oint_C P\,dx+Q\,dy=\iint_D(Q_x-P_y)dA\)。
\item
  Stokes：\(\oint_{\partial S}F\cdot d\mathbf r=\iint_S(\nabla\times F)\cdot n\,dS\)，边界方向和法向按右手规则配合。
\item
  Gauss：闭曲面外向
  \(\iint_{\partial\Omega}F\cdot n\,dS=\iiint_\Omega\nabla\cdot F\,dV\)。
\end{itemize}

\textbf{无旋、保守、无散：}\(\nabla\times F=0\) 称无旋；\(F=\nabla\phi\)
称保守，保守场的线积分路径无关，闭路积分为零。\(F\in C^1\)
在开、连通、单连通区域上无旋时保守；星形开区域是可用的充分条件。无这些区域条件可能失败，例如穿孔平面上的
\((-y/(x^2+y^2),x/(x^2+y^2))\) 无旋，但绕原点的闭路积分为
\(2\pi\)。\(\nabla\cdot F=0\)
称无散（solenoidal），与无旋是不同条件；在足够好的区域中旋度场
\(F=\nabla\times A\) 必无散，反向需要区域拓扑条件。星形区域指存在一点
\(a\)，使区域中每个点与 \(a\) 的整条线段都在区域内。

\hypertarget{cauchyriemann-ux65b9ux7a0be}{%
\subsection{65. Cauchy--Riemann
方程：E}\label{cauchyriemann-ux65b9ux7a0be}}

\textbf{原题（Original problem）}\par
Let $g$ be a differentiable function of two real variables, and let $f$ be the function of a complex variable $z$ defined by
\[f(z)=e^x\sin y+ig(x,y),\]
where $x$ and $y$ are the real and imaginary parts of $z$, respectively. If $f$ is an analytic function on the complex plane, then $g(3,2)-g(1,2)=$

(A) $e^2$ \quad (B) $e^2(\sin3-\sin1)$ \quad (C) $e^2(\cos3-\cos1)$ \quad (D) $e-e^3\sin2$ \quad (E) $(e-e^3)\cos2$.

\textbf{词语：}real and imaginary parts：实部和虚部；respectively：分别；analytic：解析的；complex plane：复平面。

\textbf{解答与分析}\par


\(u=e^x\sin y\)，\(v=g\)。解析性给出 \(u_x=v_y\)、\(u_y=-v_x\)，所以
\(g_x=-e^x\cos y\)。固定 \(y=2\)，

\[g(3,2)-g(1,2)=\int_1^3-e^x\cos2\,dx=(e-e^3)\cos2.\]

也可得 \(g=-e^x\cos y+C\)。记忆时写出矩阵关系
\(u_x=v_y,\ u_y=-v_x\)；若要从 CR
方程反推解析，通常还需要偏导连续等充分条件。

\hypertarget{ux8badux7ec3ux4e0eux8003ux573aux68c0ux67e5}{%
\subsection{训练与考场检查}\label{ux8badux7ec3ux4e0eux8003ux573aux68c0ux67e5}}

\begin{itemize}
\tightlist
\item
  第一优先：50 理想，53 积分正则性，60 量词，64 场论；另外回炉 36 开闭与
  37 投影。
\item
  易避免错误：24 齐次系统、27 模长、43 链式法则、44 代回、62 or both、65
  CR 符号。把具体检查动作写在草稿边上，而非只提醒自己``仔细''。
\item
  170 分钟约为每题 2 分 35 秒。若第一遍约 120 分钟，可预留约 50
  分钟处理卡题及复核，但需用正式环境的完整计时重复验证。推荐先标记不确定题，再按收益检查：看清题意和符号、代回、检查数量级及定义，最后处理知识缺口。
\item
  十错改善至四五错需要多挽回 5--6 题；36、37
  虽正确也可能在新卷失分。单次表现不能换算成可靠概率或分位数，需多套独立、限时试卷来判断稳定性。
\end{itemize}

\end{document}
